%% kapa: exam=2 round=37 year=2026 period=3 subject=compensation_law subject_ko=감정평가 및 보상법규 seq=1 number=1 points=40
%% source: Q-Net 감정평가사 기출문제 (cst00309) artlSeq=5268128 file=raw/1008/5268128/26년 제37회 감정평가사 2차 (3교시_감정평가 및 보상 법규).pdf
%% conversion: kordoc 4.24.51 -> tools/gen2.py (auto); review=proofread-gpt-terra; flags=none
\begin{longquestion}{【문제 1】}
국토교통부장관은 대규모의 A반도체단지가 신규 조성됨에 따라 한국철도시설공단(이하 ‘철도공단’이라 한다)을 사업시행자로 하여 고속철도를 추가 건설하는 내용의 ‘B고속철도 건설사업’(이하 ‘이 사건 사업’이라 한다)에 대한 실시계획을 승인하고 이를 고시하였다. \par
\medskip\noindent (1) 甲은 이 사건 사업구역 안에서 전기제품 제조공장(이하 ‘이 사건 공장’이라 한다)을 운영하고 있는데, 중앙토지수용위원회는 「공익사업을 위한 토지 등의 취득 및 보상에 관한 법률」(이하 ‘토지보상법’이라 한다)에 따른 적법한 절차를 거쳐 이 사건 공장의 일부를 분할 수용하는 수용재결을 하였다. \par
\medskip\noindent (2) 乙은 이 사건 사업구역 안에서 토지(이하 ‘이 사건 토지’라 한다)를 소유하고 있는데, 중앙토지수용위원회는 토지보상법에 따른 적법한 절차를 거쳐 이 사건 토지의 일부를 분할 수용하는 수용재결을 하였다. \par
\medskip\noindent (3) 丙은 이 사건 사업구역 밖에서 닭 500마리를 사육하는 양계장(이하 ‘이 사건 양계장’이라 한다)을 운영하고 있는데, 고속철도가 완공된 후 한국철도공사(이하 ‘철도공사’라 한다)가 운행하는 고속철도의 소음ㆍ진동으로 인해 이 사건 양계장의 이전이 불가피하게 되자 중앙토지수용위원회에 철도공단을 상대로 이 사건 양계장의 이전비 및 영업손실의 보상을 구하는 손실보상재결을 신청하였으나 기각되었다. \par
다음 물음에 답하시오. (40점) \par
\medskip\noindent 물음 1) 甲은 이 사건 공장의 일부 수용으로 인하여 ① ‘잔여 건축물 가격감소’ 및 ② ‘이 사건 공장의 물류이동로 변경공사비’의 손실을 입게 되었다고 주장하면서 그에 대한 추가 보상을 요구하였고, 철도공단은 甲이 추가 보상을 구하는 위 손실들에 대해서는 재결을 거치지 아니하였으므로 이는 보상의 대상이 될 수 없다고 주장한다. 한편, 중앙토지수용위원회는 위 수용재결을 할 때 ‘잔여 건축물 보수비’와 ‘이 사건 공장의 일부가 편입됨으로써 이 사건 공장 영업에 발생한 손실’이라는 보상항목에 관하여 실체적인 심리ㆍ판단을 하였으나, 甲이 추가로 주장하는 위 ① 및 ②에 관하여는 실체적인 심리ㆍ판단을 하지 않았다. 철도공단의 주장이 타당한지에 관하여 검토하시오. (15점) \par
\medskip\noindent 물음 2) 乙은 이 사건 토지 중 수용되고 남은 토지만으로는 종래의 목적대로 사용하는 것이 현저히 곤란하게 되었고, 토지의 가격도 하락하자 이에 대한 권리구제를 받고자 한다. 乙의 토지보상법상 권리구제수단에 관하여 설명하시오. (10점) \par
\medskip\noindent 물음 3) 丙이 이 사건 사업구역 밖에서 입은 손실의 특성을 설명하고 그 손실이 토지보상법상 손실보상의 대상이 되는지, 그리고 만약 손실보상의 대상이 된다고 볼 경우에는 누구를 피고로 하여 어떠한 소송으로 다투어야 하는지에 관하여 설명하시오. (15점) \par
\end{longquestion}
